\section{Каузальная окрестность одного такта}
\label{sec:epsilon}

В дискретном пространстве-времени распространение любого физического возмущения складывается из элементарных переходов за один такт времени~$\htick$.
Поскольку максимальная скорость передачи сигнала по связям решётки равна $\caSignalSpeed = \caStepSpace / \htick$ (\cref{def:step-tick}),
состояние узла~$x$ в момент $t + \htick$ может зависеть только от состояния самого узла~$x$ и тех его ближайших соседей, которые находятся на расстоянии ровно одного пространственного шага~$\caStepSpace$.

\begin{definition}[Каузальная $\varepsilon$-окрестность одного такта]
\label{def:epsilon}
Для всякого узла $x\in\carrierLattice$ его каузальной однотактовой окрестностью $N(x)\subset\carrierLattice$ называется конечное множество соседних узлов~$y$, соединённых с~$x$ светоподобными рёбрами равной длины~$\caStepSpace$, вдоль которых интервал Минковского на микроуровне обращается в нуль:
\[
  N(x) = \bigl\{\,y \in \carrierLattice \;\big|\; \|y - x\| = \caStepSpace,\quad \Delta s_{\Mlayer}^2 = \caSignalSpeed^2 \htick^2 - \|y - x\|^2 = 0\,\bigr\}.
\]
\end{definition}

Условие равенства длины всех рёбрам~$\caStepSpace$ в~\cref{def:epsilon} имеет прямой физический смысл:
любая более длинная связь (например, диагональ квадратной ячейки в окрестности Мура длины $\sqrt{2}\,\caStepSpace$ или связь со второй координационной сферой) за тот же такт~$\htick$ имела бы пространственноподобный интервал $\Delta s_{\Mlayer}^2 < 0$ и требовала бы скорости переноса выше~$\caSignalSpeed$, что запрещено условием локальной каузальности~$A_1$ (\cref{ch:axioms}).
Поэтому однотактовая окрестность~$N(x)$ исчерпывается первой координационной сферой решётки (\cref{fig:carrier-epsilon}).

\begin{figure}[htbp]
  \centering
  \input{figures/tikz/carrier-epsilon.tex}
  \caption{Каузальная $\varepsilon$-окрестность $N(x)$: множество ближайших узлов на светоподобном расстоянии~$\caStepSpace$, достижимых за один такт~$\htick$.}
  \label{fig:carrier-epsilon}
\end{figure}
